\documentclass[12pt]{ltjsarticle}

\usepackage{amsmath,amssymb,amsthm,mathtools}
\usepackage{xurl}
\usepackage[hidelinks]{hyperref}

\theoremstyle{plain}
\newtheorem{theorem}{定理}
\newtheorem{proposition}[theorem]{命題}
\newtheorem{lemma}[theorem]{補題}
\newtheorem{corollary}[theorem]{系}

\renewcommand{\proofname}{証明}

\newcommand{\val}[2]{[#1]_{#2}}
\newcommand{\NpD}{\mathcal{N}_{\mathrm{pD}}}
\newcommand{\DP}{\mathcal{D}_{\mathrm{P}}}
\newcommand{\Pset}{\bigl(\tfrac14\cdot
  \mathcal{N}_{\mathrm{pD},\,\geq12}\bigr)}
\newcommand{\seqnum}[1]{\href{https://oeis.org/#1}{\rm \underline{#1}}}

\title{原始Dyck数による加法的表示}
\author{Takayuki Kuriyama}
\date{}

\begin{document}

\maketitle

\begin{abstract}
開き括弧を $1$、閉じ括弧を $0$ としてDyck語を二進整数として読むと、
完全均衡数が得られる。本稿ではそのうち、走行平衡が末尾でのみゼロに戻る
原始Dyck語から生じる部分集合を研究する。任意の正の偶数を表示するために
必要となる原始Dyck数の最小個数を決定する。正の偶数のうち、$46$だけが
ちょうど8個の加数を必要とし、
$34$, $44$, $98$, $154$, $198$, $202$, $206$, $838$, $842$, $846$
は7個を必要とする。それ以外の正の偶数はすべて、高々6個の原始Dyck数の
和である。特に、$848$は、それ以降では常に6個の加数で十分となる鋭い
偶数閾値である。証明では、正則部分言語と、全原始族を $4$ で割った集合が
もつ自己相似的な閉包性を組み合わせる。
\end{abstract}

\noindent\textbf{キーワード：} 加法的表示、相対漸近加法基底、例外集合、
Dyck語、原始Dyck路、正則言語、区間充填。

\noindent\textbf{2020 Mathematics Subject Classification:} Primary 11B13;
Secondary 05A15, 68Q45.

\section{序論}\label{sec:introduction}

均衡のとれた括弧列はすべて、開き括弧を $1$、閉じ括弧を $0$ とすることで、
アルファベット $\{1,0\}$ 上の語として書くことができ、したがって二進整数として
読むことができる。この方法でどの整数が得られ、また得られる集合は加法的に
どれほど豊かであるのだろうか。本稿では、この問いの自然な変種を扱う。

\emph{Dyck語}とは、$1$ と $0$ を同数含み、かつ任意の接頭辞において
$1$ の個数が $0$ の個数以上であるような語 $w\in\{1,0\}^*$ をいう。
すなわち、$1$ は括弧を開き、$0$ は括弧を閉じる。空でないDyck語が
\emph{原始的}であるとは、その空でない真の接頭辞のどれも均衡していない
ことをいう。同値に、対応する格子路が終点において初めて高さゼロに戻ることをいう。
このような路は、\emph{既約}Dyck路、あるいは\emph{直既約}Dyck路とも
よく呼ばれる。原始Dyck語の言語を $\DP$ と書く。長さ順に並べ、同じ長さの
ものを辞書式順序で並べると、その最初の18個は
\[
\begin{gathered}
10,\ 1100,\ 110100,\ 111000,\ 11010100,\ 11011000,\\
11100100,\ 11101000,\ 11110000,\ 1101010100,\ 1101011000,\\
1101100100,\ 1101101000,\ 1101110000,\ 1110010100,\\
1110011000,\ 1110100100,\ 1110101000
\end{gathered}
\]
である。

$m$ 進数字上の語 $w$ に対し、それが $m$ 進法で表す整数を $\val{w}{m}$ と書く。
便宜上 $0$ を付け加えて、
\[
 \NpD=\{0\}\cup\{\val{w}{2}:w\in\DP\}
\]
とおく。最初の18個の正の元は
\[
\begin{gathered}
2,12,52,56,212,216,228,232,240,\\
852,856,868,872,880,916,920,932,936
\end{gathered}
\]
であり、これはOEIS数列 \seqnum{A057547} \cite{OEIS-A057547} である。
空でないDyck語はすべて $0$ で終わるので、$\NpD$ の正の元はすべて偶数である。
したがって、$\NpD$ の元の和によって表示できる対象整数は偶数に限られる。

以下では、次のより鋭い合同式に関する観察を繰り返し用いる。

\begin{lemma}\label{lem:mod4}
$4$ を法として $2$ と合同な正の原始Dyck数は $2$ だけである。
それ以外の正の原始Dyck数はすべて $4$ で割り切れる。
\end{lemma}

\begin{proof}
空でないDyck語はすべて $0$ で終わる。長さが4以上の原始Dyck語は、実際には
$00$ で終わらなければならない。もし $10$ で終わるなら、その路は最後の2文字を
読む直前にすでに高さゼロへ戻ってしまうからである。語 $10$ は $2$ を表す。
したがって、それより長い原始語が表す数はすべて $4$ の倍数である。
\end{proof}

Bell、Lidbetter、Shallitは、\emph{すべての}Dyck語から得られるより大きな集合、
すなわち\emph{完全均衡数}（OEIS数列 \seqnum{A014486}
\cite{OEIS-A014486}）を研究した。彼らは正則な過小近似とオートマトンを用いて、
\[
        8,\ 18,\ 28,\ 38,\ 40,\ 82,\ 166
\]
を除くすべての偶数が高々3個の完全均衡数の和である一方、漸近的にも2個の
加数では十分でないことを証明した \cite[Thm.\ 14]{BLS}。

原始語への制限は、この結果の単なる表面的な変種ではない。空でないDyck語は
すべて、原始Dyck語の連結積として一意に分解されるが、本稿では各加数そのものが
単一の因子でなければならない。したがって、制限のないDyck数による表示から、
一般には同じ個数の原始Dyck数による表示を得ることはできない。

数字によって定義される加法基底は、オートマトン理論的な方法によっても研究されてきた。
Bell、Hare、Shallitは、加法基底となる自動集合を特徴づけ、その位数を決定する
有効な手続きを与えた \cite{BHS}。関連する決定手続きにより、二進回文数
\cite{RSS} や二進平方語 \cite{MNRS} に対しても鋭い加法的結果が得られている。
原始Dyck語の言語は正則ではないため、これらの一般的な有限オートマトン法からは、
本稿の問題は直接には解決されない。

本稿の証明は、正則近似だけでは不十分となる構造的な箇所で、従来の方法と異なる。
正則部分言語によって $4$ を法として $0$ の合同類を処理し、$4$ を法として $2$ の
合同類は、全原始族の自己相似的な閉包性によって処理する。これにより、単なる
最終的な上界だけでなく、完全な例外集合、すべての偶数に対する最適な一様上界、
および鋭い最終的閾値が得られる。

長さ $2n$ のDyck語のうち、原始的なものの割合は $n\to\infty$ のとき $1/4$ に
収束する。実際、長さ $2n$ のDyck語は $C_n$ 個あり、原始Dyck語は一意に
$1u0$ と書ける。ただし $u$ は長さ $2n-2$ のDyck語である。したがって、長さ
$2n$ の原始語は $C_{n-1}$ 個である。Catalan数の公式
$C_n=\frac1{n+1}\binom{2n}{n}$ から、
\[
 \frac{C_{n-1}}{C_n}=\frac{n+1}{2(2n-1)}\longrightarrow\frac14
\]
を得る。しかし、その算術的な効果は著しい。十分大きな偶数はすべて、高々3個の
制限のないDyck数だけで表せるにもかかわらず、原始族には小さいが非自明な例外集合が
存在する。本稿では、高々6個の原始Dyck数の和でない正の偶数がちょうど11個であることを
示す。そのうち10個は7個の加数を必要とし、$46$だけが8個を必要とする。
11個の例外はすべて $848$ 未満にある。

8個の加数が必要となる理由は完全に局所的である。$48$ 未満の正の原始Dyck数は
$2$ と $12$ だけだからである。一方、最終的な6加数上界は、二つの自己相似的な族から
生じる。正則部分族が $4$ の倍数を処理し、全原始族を $4$ で割った集合が残りの
合同類を処理する。

正の偶数 $N$ に対し、$N$ を和として表すのに必要な正の原始Dyck数の最小個数を
$r(N)$ とし、本稿の中心となる整数列を
\[
        a(n)=r(2n),\qquad n\geq1
\]
によって定める。これはOEIS数列 \seqnum{A395858} \cite{OEIS-A395858} に登録された。
その最初の50項は
\[
\begin{gathered}
1,2,3,4,5,1,2,3,4,5,6,2,3,4,5,6,7,3,4,5,6,7,8,4,5,\\
1,2,1,2,3,4,2,3,2,3,4,5,3,4,3,4,5,6,4,5,4,5,6,7,5
\end{gathered}
\]
である。ここに表示した項と、以下で用いる有限証明書は、付属するPythonプログラムから
再現できる。

\begin{theorem}\label{thm:main}
$r(N)$ が6を超える場合は、正確に次のとおりである：
\[
 r(46)=8,
\]
また、
\[
 r(N)=7
 \quad\Longleftrightarrow\quad
 N\in\{34,44,98,154,198,202,206,838,842,846\}.
\]
それ以外の正の偶数はすべて $r(N)\leq6$ を満たす。したがって、任意の偶数
$N\geq848$ は、高々6個の原始Dyck数の和である。
\end{theorem}

同値な形で、数列 $a(n)=r(2n)$ を用いれば、
\[
        \max_{n\geq1}a(n)=8,
        \qquad
        a(n)=8\quad\Longleftrightarrow\quad n=23,
\]
かつ
\[
 a(n)=7
 \quad\Longleftrightarrow\quad
 n\in\{17,22,49,77,99,101,103,419,421,423\}
\]
である。したがって主定理は、単なる最終的評価ではなく、$a(n)$ の6を超える項を
すべて完全に決定している。

例外値 $46$ は、一様上界が8より小さくなり得ない理由をすでに説明している。

\begin{proposition}\label{prop:46}
整数 $46$ は、ちょうど8個の正の原始Dyck数を必要とする。すなわち、$46$ は8個の
和であるが、高々7個の和ではない。
\end{proposition}

\begin{proof}
$48$ 未満の正の原始Dyck数は $2$ と $12$ だけである。実際、長さ2と4の原始語は
それぞれ $10$ と $1100$ である。一方、長さ6以上の原始語は $11$ で始まるので、
それが表す整数は少なくとも $[110000]_2=48$ である。したがって、$46$ の任意の表示は
\[
        46=12a+2b,
        \qquad a,b\geq0
\]
の形である。同値に、$6a+b=23$ である。$12\cdot4>46$ なので $a\leq3$ であり、
したがって加数の個数は
\[
        a+b=23-5a\geq8
\]
を満たす。等号は
\[
        46=12+12+12+2+2+2+2+2
\]
によって達成される。
\end{proof}

例えば、
\[
 2026=2+240+852+932=2+216+872+936
\]
である。この二つの表示に現れる7個の相異なる整数は、いずれも原始Dyck二進展開をもつ。

以下で用いる和集合の記法は、加法基底論における標準的なものである
\cite[Chap.\ 1]{Nathanson}。非負整数 $k$ に対し、$kX$ は $X$ の元を重複を許して
$k$ 個足して得られる和の集合を表し、$0X=\{0\}$ とする。また、スカラー $c$ に対し、
$c\cdot X=\{cx:x\in X\}$ は $X$ の $c$ 倍拡大を表す。

$A\subseteq2\mathbb N_0$ に対し、$2\mathbb N_0$ の各元が $A$ の高々 $h$ 個の元の
和であるとき、$A$ を $2\mathbb N_0$ に対する\emph{位数高々 $h$ の相対加法基底}と呼ぶ。
この性質をもつ最小の整数が $h$ であるとき、その相対位数は\emph{ちょうど $h$} であるという。
同様に、$2\mathbb N_0$ の十分大きな元がすべて $A$ の高々 $h$ 個の元の和であるとき、
$A$ を\emph{位数高々 $h$ の相対漸近加法基底}と呼ぶ。

\section{原始語の内部にある正則族}

原始Dyck語の中には、単純な正則パターンが通っている。例えば、$868$ の二進展開は
\[
 \operatorname{bin}(868)=1101100100=11\,01\,10\,01\,00
 \in 11(01\mid10)^*00
\]
と因数分解される。二進数字を2個ずつ組にすると、最初のブロック $11$ は4進数字 $3$、
最後のブロック $00$ は $0$ となり、中間の各ブロック $01$ または $10$ は $1$ または
$2$ となる。このことは、このような4進整数を固定個数足した和が長い区間を埋める可能性を
示唆する。以下の混合和集合補題によって、この考えを厳密にする。

\[
 \mathcal E=\{0\}\cup
 \bigl\{\val{3\varepsilon_1\cdots\varepsilon_k}{4}:
        k\geq0,\ \varepsilon_i\in\{1,2\}\bigr\}
\]
とおく。$k=0$ のとき、文字列 $\varepsilon_1\cdots\varepsilon_k$ は空であるため、
対応する元は $\val{3}{4}=3$ である。$4$ 倍することは、4進表示の末尾に $0$ を付ける
ことに対応する。各4進数字を対応する二進数字2個で置き換えると、
\[
\begin{aligned}
 3_{(4)}&\longleftrightarrow11_{(2)},&
 1_{(4)}&\longleftrightarrow01_{(2)},\\
 2_{(4)}&\longleftrightarrow10_{(2)},&
 0_{(4)}&\longleftrightarrow00_{(2)}
\end{aligned}
\]
である。したがって、$4\cdot\mathcal E$ の正の元の二進展開はすべて
$11(01\mid10)^*00$ に属する。最初の $11$ を読んだ後、路は高さ $2$ にあり、
中間の各ブロック $01$ または $10$ は路を高さ $2$ に戻し、最後の $00$ が初めて
高さゼロへ戻す。ゆえに、$4\cdot\mathcal E$ の正の元はすべて原始Dyck数である。
また $0\in\NpD$ なので、
\begin{equation}\label{eq:4E}
        4\cdot\mathcal E\subseteq\NpD
\end{equation}
を得る。

核心は、$\mathcal E$ の6個の元だけで、$25$ 以降のすべての整数を覆えることである。

\begin{proposition}\label{prop:E}
任意の整数 $n\geq25$ は $6\mathcal E$ に属する。
\end{proposition}

これを、有限桁に基づく区間論法によって証明する。まず、以下の補助集合の役割を説明して
おく価値がある。$\mathcal E_m$ の $m$ 桁の元を先頭桁と下位桁に分けると、下位
$m-1$ 桁は、各桁が $\{0,1\}$ に属する4進補正項をなす。集合 $\Delta_m$ は、まさに
これらの補正項を記録する。6項和の中のいくつかの加数が新しい先頭桁を用いるとき、残る
下位桁の問題は、$\Delta_{m-1}$ と $\mathcal E_{m-1}$ のコピーの混合和となる。
混合和集合 $S_{t,m}$ は、これらすべての可能性を同時に追跡するために導入される。

$u_0=0$ とおき、$m\geq1$ に対して
\begin{equation}\label{eq:umdef}
 u_m=[\underbrace{11\cdots1}_{m\text{ 桁}}]_4
    =1+4+\cdots+4^{m-1}=\frac{4^m-1}{3}
\end{equation}
と定める。このとき
\begin{equation}\label{eq:urecursion}
\begin{aligned}
 4^{m-1}&=3u_{m-1}+1 &&(m\geq1),\\
 u_{m-1}&=4u_{m-2}+1 &&(m\geq2)
\end{aligned}
\end{equation}
が成り立つ。補正集合を
\[
 \Delta_m=\Bigl\{[\delta_{m-1}\cdots\delta_0]_4:\delta_j\in\{0,1\}\Bigr\}
    =\Bigl\{\textstyle\sum_{j=0}^{m-1}\delta_j4^j:\delta_j\in\{0,1\}\Bigr\}
\]
と定める。このとき $\max\Delta_m=u_m$ である。$\mathcal E_m$ を、$0$ と、4進表示が
高々 $m$ 桁である $\mathcal E$ の元とからなる有限部分集合とする。例えば、
\[
\begin{aligned}
 \Delta_2&=\{[00]_4,[01]_4,[10]_4,[11]_4\}=\{0,1,4,5\},\\
 \mathcal E_2&=\{0,[3]_4,[31]_4,[32]_4\}=\{0,3,13,14\}
\end{aligned}
\]
である。$0\leq t\leq6$ に対して
\[
        S_{t,m}=t\Delta_m+(6-t)\mathcal E_m
\]
とおく。$\mathcal E_m$ の最大元は
\[
 \max\mathcal E_m=[3\underbrace{22\cdots2}_{m-1\text{ 桁}}]_4
                =3\cdot4^{m-1}+2u_{m-1}
\]
であり、$m\geq2$ では式 \eqref{eq:urecursion} から
\begin{equation}\label{eq:maxErecursion}
 \max\mathcal E_{m-1}=11u_{m-2}+3=\frac{11u_{m-1}+1}{4}
\end{equation}
を得る。$\max\Delta_m=u_m$ なので、
\begin{equation}\label{eq:maxStm}
 \max S_{t,m}=tu_m+(6-t)\max\mathcal E_m
\end{equation}
である。さらに $\max\mathcal E_m>u_m$ だから、$\max S_{t,m}$ は $t$ とともに減少する。
以下、$[a,b]$ は整数区間 $\{a,a+1,\ldots,b\}$ を表す。

\begin{lemma}\label{lem:mixed}
任意の $m\geq2$ と任意の $2\leq t\leq6$ に対して、
\begin{equation}\label{eq:mixedtge}
 S_{t,m}=[0,\max S_{t,m}]
\end{equation}
が成り立つ。一方、
\begin{equation}\label{eq:mixedtone}
 S_{1,m}=[0,\max S_{1,m}]\setminus\{2\}
\end{equation}
である。さらに、
\begin{equation}\label{eq:mixedEminterval}
 [25,\ell_m]\subseteq6\mathcal E_m,
 \qquad \ell_m=\frac{33}{8}\,4^m-4
\end{equation}
が成り立つ。
\end{lemma}

\begin{proof}
まず基底の場合 $m=2$ において、\eqref{eq:mixedtge} と
\eqref{eq:mixedtone} を証明する。すでに
\[
 \Delta_2=\{0,1,4,5\},\qquad \mathcal E_2=\{0,3,13,14\}
\]
であることを見た。直接計算すると
$2\Delta_2=\{0,1,2,4,5,6,8,9,10\}$ であり、さらにもう1個
$\Delta_2$ を加えると残る隙間がすべて埋まる。したがって
$3\Delta_2=[0,15]$ であり、よって
\[
 t\Delta_2=[0,5t]\qquad(3\leq t\leq6)
\]
である。一方、
\[
\begin{aligned}
2\mathcal E_2
={}&\{0,3,6,13,14,16,17,26,27,28\},\\[1ex]
3\mathcal E_2
={}&\{0,3,6,9\}\cup[13,14]\cup[16,17]\cup[19,20]\\
&\qquad\cup[26,31]\cup[39,42],\\[1ex]
4\mathcal E_2
={}&\{0,3,6,9\}\cup[12,14]\cup[16,17]\cup[19,20]\\
&\qquad\cup[22,23]\cup[26,34]\cup[39,45]\cup[52,56],\\[1ex]
5\mathcal E_2
={}&\{0,3,6,9\}\cup[12,17]\cup[19,20]\cup[22,23]\\
&\qquad\cup[25,37]\cup[39,48]\cup[52,59]\cup[65,70]
\end{aligned}
\]
である。$5\mathcal E_2$ に $\Delta_2$ を加え、$4\mathcal E_2$ に
$2\Delta_2$ を加えると、
\[
 S_{1,2}=\{0,1\}\cup[3,75],\qquad S_{2,2}=[0,66]
\]
を得る。同様に、
\[
\begin{aligned}
S_{3,2}&=3\Delta_2+3\mathcal E_2=[0,15]+3\mathcal E_2=[0,57],\\[1ex]
S_{4,2}&=4\Delta_2+2\mathcal E_2=[0,20]+2\mathcal E_2=[0,48],\\[1ex]
S_{5,2}&=5\Delta_2+\mathcal E_2=[0,25]+\mathcal E_2=[0,39],\\[1ex]
S_{6,2}&=6\Delta_2=[0,30]
\end{aligned}
\]
である。以上をまとめると、
\[
\begin{array}{c|c}
t&S_{t,2}\\ \hline
1&[0,75]\setminus\{2\}\\
2&[0,66]\\
3&[0,57]\\
4&[0,48]\\
5&[0,39]\\
6&[0,30]
\end{array}
\]
となり、$m=2$ に対する \eqref{eq:mixedtge} と
\eqref{eq:mixedtone} が証明された。

\eqref{eq:mixedEminterval} については、
$[25,\ell_2]=[25,62]\subseteq6\mathcal E_2$ を確認すればよい。
区間 $[26,28]$, $[39,42]$, $[52,56]$ は、それぞれ
$\{13,14\}$ の2個、3個、4個の元の和であるから、
\[
\begin{aligned}
25&=13+3+3+3+3+0,\\
[26,40]&=\{26,27,28\}+3\cdot\{0,1,2,3,4\},\\
[39,51]&=[39,42]+3\cdot\{0,1,2,3\},\\
[52,62]&=[52,56]+3\cdot\{0,1,2\}
\end{aligned}
\]
である。これで基底の場合が完了する。下端 $25$ は鋭い。実際、$25$ 未満の
$\mathcal E$ の正の元は $3,13,14$ だけであり、これらの高々6個の和として
$24$ を表すことはできないので、$24\notin6\mathcal E$ である。

次に $m\geq3$ とし、レベル $m-1$ で
\eqref{eq:mixedtge}--\eqref{eq:mixedEminterval} が成り立つと仮定する。
この段階での記法を簡潔にするため、
\[
 u=u_{m-1},\qquad q=q_m,
 \qquad M=\max\mathcal E_{m-1}
\]
と書く。このとき
\[
 4^{m-1}=3u+1,\qquad q=10u+3,
 \qquad M=\frac{11u+1}{4}
\]
である。ここで
\begin{equation}\label{eq:qmdef}
 q=q_m=[3\underbrace{11\cdots1}_{m-1\text{ 桁}}]_4
      =3\cdot4^{m-1}+u=10u+3
\end{equation}
であり、以下では
\begin{equation}\label{eq:qoverlap}
 2\cdot4^{m-1}+4u=q-1
\end{equation}
を繰り返し用いる。4進表示の先頭桁によって分けると、
\[
 \Delta_m=\Delta_{m-1}\cup(4^{m-1}+\Delta_{m-1}),\qquad
 \mathcal E_m=\mathcal E_{m-1}\cup(q+\Delta_{m-1})
\]
を得る。実際、$\Delta_m$ の元の新しい先頭桁は $0$ または $1$ である。一方、
$\mathcal E_m$ の新しい $m$ 桁の元は、先頭桁が $3$ で、その後に
$\Delta_{m-1}$ から来る下位 $m-1$ 桁の補正項が続く。

$t$ 個の $\Delta_m$ のコピーのうち $s$ 個が平行移動された部分を用い、
$6-t$ 個の $\mathcal E_m$ のコピーのうち $v$ 個が新しい $m$ 桁部分を用いるなら、
\begin{equation}\label{eq:decomp}
 S_{t,m}=\bigcup_{v=0}^{6-t}\ \bigcup_{s=0}^{t}
 \bigl(s\cdot4^{m-1}+vq+S_{t+v,m-1}\bigr)
\end{equation}
である。したがって、$v$ を固定すると、パラメータ $s$ は
$4^{m-1}$ の倍数だけ平行移動した $t+1$ 個の集合を生む。さらに $v$ を1増やすと、
次の大きなブロックへ移る：
\[
 \underbrace{[vq,\ \cdots]}_{v\text{-ブロック}}
 \qquad
 \underbrace{[(v+1)q,\ \cdots]}_{(v+1)\text{-ブロック}}.
\]
以下では、隣り合う平行移動区間と、隣り合う $v$-ブロックが重なることを確認する。

まず $t\geq2$ の場合を扱う。帰納法の仮定 \eqref{eq:mixedtge} により、
\eqref{eq:decomp} に現れるすべての $v$ に対して
$S_{t+v,m-1}=[0,\max S_{t+v,m-1}]$ である。$v$ を固定したとき、
連続する $s$ の値はこの区間を $4^{m-1}$ だけずつ平行移動する。
$\max S_{j,m-1}$ は $j$ とともに減少するので、
\[
 \max S_{t+v,m-1}\geq\max S_{6,m-1}=6u
     \geq3u=4^{m-1}-1
\]
である。したがって、これら $t+1$ 個の平行移動区間は互いに重なり、
\[
 [vq,\,vq+t\cdot4^{m-1}+\max S_{t+v,m-1}]
\]
をなす。

連続する $v$ に対応するブロック同士も重なる。次のブロックが存在するときは
$t+v\leq5$ であり、また $t\geq2$ だから、
\[
\begin{aligned}
 t\cdot4^{m-1}+\max S_{t+v,m-1}
 &\geq2\cdot4^{m-1}+\max S_{5,m-1}\\
 &=2\cdot4^{m-1}+5u+M\\
 &=(2\cdot4^{m-1}+4u)+(u+M)\\
 &\geq q-1
\end{aligned}
\]
が \eqref{eq:qoverlap} により成り立つ。ゆえに、\eqref{eq:decomp} の和集合は、
$0$ からその最大元 $\max S_{t,m}$ までの全整数区間である。これでレベル $m$ における
\eqref{eq:mixedtge} が証明された。

$t=1$ の場合、最初のブロックには欠けた整数 $2$ が一つだけ残る。固定した
$v\geq1$ に対し、二つの $s$-平行移動は $S_{1+v,m-1}$ を含み、$1+v\geq2$ である。
また、
\[
 \max S_{1+v,m-1}\geq\max S_{6,m-1}=6u\geq4^{m-1}-1
\]
なので、この二つの平行移動区間は重なる。したがって、それ以降の各 $v$-ブロックは
区間である。$\max S_{j,m-1}$ は $j$ とともに減少するため、連続する
$v$-ブロックの最後の組だけを確認すれば十分である。実際、
\[
\begin{aligned}
 4^{m-1}+\max S_{5,m-1}-(q-1)
 &=M-2u-1\\
 &=\frac{11u+1}{4}-2u-1
  =\frac{3u-3}{4}\geq0
\end{aligned}
\]
なので、連続する $v$-ブロックは重なる。最後に、
\[
\begin{aligned}
 \max S_{1,m-1}-(4^{m-1}+2)
 &=u+5M-4^{m-1}-2\\
 &=\frac{47u-7}{4}\geq0
\end{aligned}
\]
である。したがって、最初の $v$-ブロック内で、平行移動していない長い区間は
$4^{m-1}+2$ まで達し、次の $s$-平行移動区間と隙間なく接続する。ゆえに、
\[
 S_{1,m}=[0,\max S_{1,m}]\setminus\{2\}
\]
を得る。これでレベル $m$ における \eqref{eq:mixedtone} が証明された。

残るのは、$6\mathcal E_m$ 内の区間を次のレベルへ伝播させることである。
\eqref{eq:decomp} で $t=0$ とすると、
\[
 6\mathcal E_m=\bigcup_{v=0}^{6}\bigl(vq+S_{v,m-1}\bigr)
\]
である。$v=0$ のブロックは $S_{0,m-1}=6\mathcal E_{m-1}$ であり、
帰納法の仮定により $[25,\ell_{m-1}]$ を含む。$m\geq3$ なので $u\geq5$ であり、
\[
 \ell_{m-1}-(q+2)=\frac{19u-39}{8}\geq0
\]
である。$v=1$ のブロック、すなわち $q+S_{1,m-1}$ は、帰納法の仮定により
$\{q,q+1\}\cup[q+3,q+\max S_{1,m-1}]$ を含む。直前の区間は $q+2$ まで
達しているので、この二つは隙間なく接続する。

$v=j$ のブロックを $v=j+1$ のブロックへ接続するためには、$j=1,2,3$ について
\[
 \max S_{j,m-1}\geq q-1
\]
であれば十分である。$\max S_{j,m-1}$ は $j$ とともに減少するので、最も厳しいのは
$j=3$ の場合である。したがって、
\[
\begin{aligned}
 \max S_{3,m-1}-(q-1)
 &=3u+3M-(10u+2)\\
 &=\frac{5u-5}{4}\geq0
\end{aligned}
\]
を確認すればよい。ゆえに
$\max S_{1,m-1},\ \max S_{2,m-1},\ \max S_{3,m-1}\geq q-1$ であり、
$v=1,2,3,4$ のブロックは順に接続して、
\[
 4q+\max S_{4,m-1}=\frac{33}{8}\,4^m-4=\ell_m
\]
までを覆う。残る $v=5,6$ のブロックは隙間の後に始まる可能性があるが、ここで伝播させる
区間には不要である。これでレベル $m$ における
\eqref{eq:mixedtge}--\eqref{eq:mixedEminterval} が証明され、帰納法が完了する。
\end{proof}

$\ell_m\to\infty$ であり、$\mathcal E=\bigcup_{m\geq1}\mathcal E_m$ だから、
命題 \ref{prop:E} は直ちに従う。

前の命題によって、$4$ を法として $0$ の合同類を処理できる。すなわち、$100$ 以上の
任意の $4$ の倍数は、高々6個の原始Dyck数の和である。

\section{全原始族を4で割った集合}

正則族 $\mathcal E$ は $4$ の倍数を扱うには十分であるが、以下で必要となる5加数区間を
それだけで与えることはできない。そこで、もう一方の合同類に対しては全原始族を用いる。
\[
 \Pset=\{v/4:v\in\NpD,\ v\geq12\}
\]
とおく。したがって、
\[
 \NpD=\{0,2\}\cup(4\cdot\Pset)
\]
である。整数集合 $X$ と $k\geq0$ に対して、
\[
 F_k(X)=\bigcup_{j=0}^{k}jX
\]
と書き、これを $X$ の高々 $k$ 個の元の和全体の集合とする。

全原始族は、単純な自己相似的閉包性をもつ。

\begin{lemma}\label{lem:Pclosure}
$p\in\Pset$ ならば、$4p+1$ および $4p+2$ も $\Pset$ に属する。
\end{lemma}

\begin{proof}
原始Dyck数 $4p$ の二進展開を $u00$ と書く。$u$ を読み終えた直後、対応する路の高さは
$2$ である。$4(4p+1)$ と $4(4p+2)$ の二進展開は、それぞれ
\[
 u0100\qquad\text{および}\qquad u1000
\]
である。挿入されたブロック $01$ と $10$ は、いずれも高さ $2$ から始まり高さ $2$ で
終わり、その途中の高さも正である。したがって、どちらの新しい路も最後の記号まで
厳密にゼロより上にとどまり、末尾でのみゼロへ戻る。ゆえに、表示された二つの語は
原始Dyck語であり、主張が従う。
\end{proof}

次に、$\Pset$ に対する5加数区間を確立する。有限の基底部分では、二つの部分集合
\[
 L=\{3,13,14,53,54,57,58,60\}\subseteq\Pset
\]
および
\[
\begin{split}
 H=\{&213,214,217,218,220,229,230,233,234,236,\\
     &241,242,244,248\}\subseteq\Pset
\end{split}
\]
を用いる。次の表は、対応する原始Dyck数をすべて表示している。各項において、二進語
$[4h]_2$ は原始的であり、したがって $h\in\Pset$ である。
\[
\begin{array}{c|c@{\qquad}c|c}
\multicolumn{4}{c}{4L}\\
h&[4h]_2&h&[4h]_2\\ \hline
3&1100&53&11010100\\
13&110100&54&11011000\\
14&111000&57&11100100\\
58&11101000&60&11110000
\end{array}
\]
\[
\begin{array}{c|c@{\qquad}c|c}
\multicolumn{4}{c}{4H}\\
h&[4h]_2&h&[4h]_2\\ \hline
213&1101010100&233&1110100100\\
214&1101011000&234&1110101000\\
217&1101100100&236&1110110000\\
218&1101101000&241&1111000100\\
220&1101110000&242&1111001000\\
229&1110010100&244&1111010000\\
230&1110011000&248&1111100000
\end{array}
\]

有限計算は、コミット
\path{bc0211ccb478d79969f8ce6b5fe465e5f3fd9cf5}
に保存されたPythonプログラム
\texttt{verify\_certificates\_for\_paper.py} を用いて再現できる：
\url{https://github.com/growupkuriyama-hub/lean_cfg_project/blob/bc0211ccb478d79969f8ce6b5fe465e5f3fd9cf5/LeanCfgProject/PrimDyck/verify_certificates_for_paper.py}。
同じプログラムを補助ファイルとしても添付し、実行コマンドと期待される出力を記した
READMEも同梱する。

補題 \ref{lem:finitefive} の五つの主張は包含証明書である。すなわち、プログラムは
各表示区間が対応する和集合に含まれることを確認するのであって、和集合がその区間に
等しいと主張するものではない。これに対し、補題 \ref{lem:finiteexception} については、
プログラムは $[0,211]$ に切り詰めた関連集合を正確に計算し、表示された区間の和集合との
等号を確認する。

\begin{lemma}[有限区間証明書]\label{lem:finitefive}
次の包含関係が成り立つ：
\[
\begin{array}{c|c}
\text{和集合}&\text{含まれる整数区間}\\ \hline
5L &[215,254]\\
H+4L &[245,486]\\
2H+3L &[439,674]\\
3H+2L &[645,862]\\
4H+L &[855,1046].
\end{array}
\]
\end{lemma}

\begin{proof}
\[
 A=\{53,54,57,58,60\},\qquad B=\{3,13,14\}
\]
とおく。逐次加算により、
\[
 2A=[106,108]\cup[110,118]\cup\{120\},
\]
\[
 3A=[159,178]\cup\{180\},\qquad
 4A=[212,238]\cup\{240\}
\]
を得る。したがって $4A+B=[215,254]$ であり、第1行が証明される。

第2行について、$L$ を繰り返し加えると、
\[
\begin{split}
4L\supseteq{}&[32,34]\cup[42,45]\cup[52,56]\cup[82,102]\\
 &{}\cup[116,148]\cup[162,194]\cup[212,238]
\end{split}
\]
を得る。最初の三つの短い区間から、
\[
\begin{aligned}
H+[32,34]&\supseteq[245,254]\cup[261,270]
 \cup[273,278]\cup[280,282],\\
H+[42,45]&\supseteq[255,265]\cup[271,281]\cup[283,293],\\
H+[52,56]&\supseteq[265,276]\cup[281,304]
\end{aligned}
\]
を得る。これらの和集合は $[245,304]$ を含む。$H$ の連続する元の差は高々 $9$ である。
したがって、$I=[a,b]$ が $b-a\geq8$ を満たすなら、
$H+I=[213+a,248+b]$ である。これを残りの四つの長い区間に適用すると、
\[
 [295,350],\quad[329,396],\quad[375,442],\quad[425,486]
\]
を得る。これらは互いに重なるので、$[245,486]\subseteq H+4L$ の証明が完了する。

第3行については、次の包含区間で十分である：
\[
 2H\supseteq[426,428]\cup[430,438]\cup[446,478]\cup\{496\},
\]
\[
\begin{split}
3L\supseteq{}&\{9\}\cup[19,20]\cup[39,42]\cup[73,77]
 \cup[109,111]\\
 &{}\cup[113,134]\cup[159,178]\cup\{180\}.
\end{split}
\]
$2H$ の最初の二つの成分から、
\[
 [430,438]+9=[439,447],\quad
 [426,428]+[19,20]=[445,448]
\]
および
\[
 [430,438]+[19,20]=[449,458]
\]
を得る。$[446,478]$ に、順に
\[
\begin{gathered}
 \{9\},\ [39,42],\ [73,77],\ [109,111],\\
 [113,134],\ [159,178],\ \{180\}
\end{gathered}
\]
を加えると、互いに重なって $[455,658]$ を覆う区間が得られる。最後に、
\[
 496+[159,178]=[655,674]
\]
である。以上の区間を合わせると、$[439,674]\subseteq2H+3L$ が証明される。

第4行については、
\[
 3H\supseteq[639,734]\cup\{744\},\qquad
 2L\supseteq\{6\}\cup[70,74]\cup[110,118]\cup\{120\}
\]
を用いる。四つの区間和
\[
 [645,740],\quad[709,808],\quad[759,854],\quad[854,862]
\]
は $[645,862]$ を覆う。最後の行については、
\[
 4H\supseteq[852,982]\cup[984,986],
 \qquad \{3,57,58,60\}\subseteq L
\]
を用いる。このとき、区間
\[
 [855,985],\quad[912,1042],\quad[1041,1044],\quad[1044,1046]
\]
が $[855,1046]$ を覆う。
\end{proof}

補題 \ref{lem:finitefive} の五つの区間は互いに重なるので、
\[
 [215,1046]\subseteq5\Pset
\]
を得る。

\begin{proposition}\label{prop:fiveP}
任意の整数 $n\geq215$ は、$\Pset$ のちょうど5個の元の和である。したがって、
\[
 [212,\infty)\subseteq F_5(\Pset)
\]
である。下端は鋭い。すなわち、$209,210,211$ のいずれも
$F_5(\Pset)$ に属さない。
\end{proposition}

\begin{proof}
補題 \ref{lem:finitefive} は、より強い包含
$[215,1046]\subseteq5\Pset$ を与える。$n$ に関する強い帰納法の基底として必要なのは、
そのうち最初の区間 $[215,867]$ だけである。$n\geq868$ とし、
$S\equiv n\pmod4$ を満たす一意な $S\in\{5,6,7,8\}$ を選ぶ。このとき、
\[
 M=\frac{n-S}{4}\geq\frac{868-8}{4}=215
\]
である。$M<n$ なので、帰納法の仮定により
\[
 M=q_1+\cdots+q_5,\qquad q_i\in\Pset
\]
と書ける。$5\leq S\leq8$ であるから、
$e_1+\cdots+e_5=S$ となる $e_i\in\{1,2\}$ を選ぶことができる。したがって、
\[
 n=4M+S=\sum_{i=1}^{5}(4q_i+e_i)
\]
である。補題 \ref{lem:Pclosure} により、各加数 $4q_i+e_i$ は $\Pset$ に属する。
ゆえに $n\in5\Pset$ であり、帰納法が完了する。

最後に、
\[
\begin{aligned}
212&=53+53+53+53,\\
213&=53+53+53+54,\\
214&=53+53+54+54
\end{aligned}
\]
であるから、$[212,\infty)\subseteq F_5(\Pset)$ である。
次に、$212$ 未満の $\Pset$ の元が、ちょうど $L$ の8元であることを確認する。
長さ $4$, $6$, $8$ の原始Dyck語を $4$ で割って得られる集合は、それぞれ
\[
 \{3\},\qquad \{13,14\},\qquad \{53,54,57,58,60\}
\]
である。長さ $10$ の原始Dyck語はすべて $1u0$ の形であり、ここで $u$ は長さ $8$ の
Dyck語である。辞書式順序で最小のそのような語は $1101010100$ であり、その値は
$852$ である。$4$ で割ると $213$ になる。さらに長い原始語は、これより大きな値をもつ。
したがって、
\[
 \Pset\cap[0,211]=L
\]
である。

$L$ の高々5個の元の和で、$\{53,54,57,58,60\}$ から高々3項を用いるものは、
最大でも
\[
 3\cdot60+2\cdot14=208
\]
である。一方、この集合から少なくとも4項を用いる和は、最小でも
\[
 4\cdot53=212
\]
である。したがって、$209,210,211\notin F_5(\Pset)$ である。
\end{proof}

\section{例外集合}\label{sec:exceptions}

まず、$4$ を法として $2$ の合同類における6加数表示可能性を、$\Pset$ 上の有限和条件へ
翻訳する。

\begin{lemma}\label{lem:sixcriterion}
$N=4m+2$ とする。このとき、$N$ が高々6個の正の原始Dyck数の和であるための必要十分条件は、
\[
 m\in F_5(\Pset)
 \cup\bigl(1+F_3(\Pset)\bigr)
 \cup\bigl(2+F_1(\Pset)\bigr)
\]
である。
\end{lemma}

\begin{proof}
補題 \ref{lem:mod4} により、$2$ 以外の原始Dyck加数はすべて $4$ の倍数である。
したがって、$N\equiv2\pmod4$ の表示には、奇数個の $2$ が含まれる。
加数の総数が高々6なら、その個数は $1$, $3$, $5$ のいずれかである。

$2$ が1個の場合、残る和を $4$ で割ると $m\in F_5(\Pset)$ を得る。
$2$ が3個の場合は $m-1\in F_3(\Pset)$ を得て、5個の場合は
$m-2\in F_1(\Pset)$ を得る。これで必要性が示された。逆に、この三つの所属条件の
いずれかが成り立てば、対応する $\Pset$ の加数を $4$ 倍し、順に1個、3個、5個の $2$ を
付け加えることで、$N$ の表示が得られる。したがって、この条件は十分でもある。
\end{proof}

\begin{lemma}[小さい4の倍数]\label{lem:smallmultiples}
$100$ 未満の正の $4$ の倍数のうち、高々6個の正の原始Dyck数で表示できないものは
$44$ だけである。さらに、
\[
 44=12+12+12+2+2+2+2
\]
である。
\end{lemma}

\begin{proof}
$100$ 未満の正の原始Dyck数は、正確に
\[
 2,\ 12,\ 52,\ 56
\]
である。まず $N<52$ とし、$N=4x$ と書く。$12$ のコピーと、対 $2+2$ の和は
\[
 x=3a+c
\]
の形であり、$a+2c$ 個の加数を用いる。条件 $a+2c\leq6$ のもとで得られる
$x<13$ の正の値は
\[
 1,2,3,4,5,6,7,8,9,10,12
\]
である。欠ける値は $x=11$ だけであり、これは $N=44$ に対応する。

次に $52\leq N<100$ とする。まず $52=4\cdot13$ または $56=4\cdot14$ のどちらかを
1個用いる。残る高々5個の加数として $12$ と対 $2+2$ を用いて得られる剰余商は
\[
 R=\{3a+c:a,c\geq0,\ a+2c\leq5\}
   =\{0,1,2,3,4,5,6,7,9,10,12,15\}
\]
である。この12個の整数を直接確認すると、
\[
 [13,24]\subseteq(13+R)\cup(14+R)
\]
である。したがって、$52$ から $96$ までのすべての $4$ の倍数は所要の表示をもつ。
最後に、上に表示した $44$ の7項表示によって証明が完了する。
\end{proof}

\begin{lemma}[有限例外集合証明書]\label{lem:finiteexception}
\[
 L=\{3,13,14,53,54,57,58,60\}
\]
とする。このとき、
\[
\begin{aligned}
F_5(L)\cap[0,211]
={}&\{0,3,6,9\}\cup[12,17]\cup[19,20]\cup[22,23]\\
 &{}\cup[25,37]\cup[39,48]\cup[52,208]
\end{aligned}
\]
である。さらに、
\[
\begin{aligned}
[0,211]\cap\Bigl(F_5(L)&\cup(1+F_3(L))\cup(2+F_1(L))\Bigr)\\
={}&[0,7]\cup[9,10]\cup[12,23]\cup[25,37]\\
 &{}\cup[39,48]\cup[52,208]
\end{aligned}
\]
である。したがって、$[0,211]$ における補集合は
\[
 \{8,11,24,38,49,50,51,209,210,211\}
\]
である。
\end{lemma}

\begin{proof}
$0L=\{0\}$ から始め、漸化式 $jL=(j-1)L+L$ を用いる。
$L$ の8個の元を順次加えると、$211$ 以下では
\[
\begin{aligned}
F_5(L)\cap[0,211]={}&\{0,3,6,9\}\cup[12,17]\cup[19,20]\cup[22,23]\\
 &{}\cup[25,37]\cup[39,48]\cup[52,208]
\end{aligned}
\]
を得る。一方、
\[
\begin{aligned}
(1+F_3(L))\cap[0,211]
={}&\{1,4,7,10\}\cup[14,15]\cup[17,18]\cup[20,21]\\
 &{}\cup[27,32]\cup[40,43]\cup[54,55]\cup[57,62]\\
 &{}\cup[64,65]\cup[67,78]\cup[80,89]\cup[107,135]\\
 &{}\cup[160,179]\cup\{181\}
\end{aligned}
\]
であり、また
\[
 (2+F_1(L))\cap[0,211]
 =\{2,5\}\cup[15,16]\cup[55,56]\cup[59,60]\cup\{62\}
\]
である。表示された区間リストの和集合を取ると第2の式を得る。そこから欠ける整数は、
最後に表示した集合の元と正確に一致する。各段階では、直前の区間リストを $L$ の8元だけ
平行移動するだけなので、この計算は全探索を行わなくても検証できる。
\end{proof}

\begin{proof}[定理 \ref{thm:main} の証明]
まず $N\equiv0\pmod4$ とする。$N\geq100$ なら、$N=4n$ と書く。このとき
$n\geq25$ なので、命題 \ref{prop:E} と \eqref{eq:4E} により、$N$ は高々6個の
原始Dyck数の和として表せる。補題 \ref{lem:smallmultiples} は残る正の $4$ の倍数を
処理し、$r(44)=7$ を示す。

次に $N\equiv2\pmod4$ とし、$N=4m+2$ と書く。$N\geq850$ なら $m\geq212$ なので、
命題 \ref{prop:fiveP} により $m\in F_5(\Pset)$ である。補題
\ref{lem:sixcriterion} により、$N$ は高々6個の原始Dyck数の和である。
前段落と合わせると、すでに、任意の偶数 $N\geq848$ がそのような表示をもつことが分かる。

残るのは、$N\leq846$ かつ $N\equiv2\pmod4$ である数を分類することである。
この場合 $0\leq m\leq211$ であり、命題 \ref{prop:fiveP} の証明で示したように、
$m$ 以下の $\Pset$ の元は
\[
 L=\{3,13,14,53,54,57,58,60\}
\]
に属する。補題 \ref{lem:finiteexception} により、表示条件で覆われない $m$ の値は
\[
 \{8,11,24,38,49,50,51,209,210,211\}
\]
である。補題 \ref{lem:sixcriterion} により、対応する $N=4m+2$ は正確に
\[
 34,46,98,154,198,202,206,838,842,846
\]
である。したがって、これら10個と $44$ を合わせたものが、高々6個の原始Dyck数の和でない
正の偶数の全体である。

これらのうち $46$ を除くすべては、次の7項表示をもつ：
\[
\begin{aligned}
34&=12+12+2+2+2+2+2,\\
44&=12+12+12+2+2+2+2,\\
98&=56+12+12+12+2+2+2,\\
154&=52+52+12+12+12+12+2,\\
198&=56+52+52+12+12+12+2,\\
202&=56+56+52+12+12+12+2,\\
206&=56+56+56+12+12+12+2,\\
838&=240+240+240+52+52+12+2,\\
842&=240+240+240+56+52+12+2,\\
846&=240+240+240+56+56+12+2.
\end{aligned}
\]
命題 \ref{prop:46} により、$46$ はちょうど8個の加数を必要とする。
これで定理のすべての主張が証明された。
\end{proof}

\begin{corollary}\label{cor:threshold}
$848$ は、任意の偶数 $N\geq T$ が高々6個の原始Dyck数の和となるような偶数 $T$ のうち、
最小のものである。
\end{corollary}

\begin{proof}
定理 \ref{thm:main} により、$848$ 以上の任意の偶数には6個の加数で十分である。
一方、$846$ は7個の加数を必要とする。
\end{proof}

\begin{corollary}\label{cor:eight}
任意の非負偶数は $\NpD$ の高々8個の元の和であり、$46$ は $\NpD$ の高々7個の元の和でない
唯一の偶数である。
\end{corollary}

\begin{proof}
正の整数に対する主張は定理 \ref{thm:main} から直ちに従う。整数 $0$ は空和である。
\end{proof}

定理 \ref{thm:main} は、6加数表示可能性に対する完全な例外集合を与える。
系 \ref{cor:eight} により、$\NpD$ は $2\mathbb N_0$ に対する相対加法基底であり、
その正確な位数は8である。また、これは位数高々6の相対漸近加法基底でもある。
同値に、$\frac12\cdot\NpD$ は $\mathbb N_0$ の位数高々6の漸近加法基底である。
系 \ref{cor:threshold} は、対応する閾値 $848$ が鋭いことを示す。

十分大きな任意の偶数を高々 $h$ 個の原始Dyck数の和として表せるような最小の $h$ を
$h_{\mathrm{asym}}$ とする。すべての完全均衡数からなるより大きな族は、位数2の漸近加法基底ではない
\cite[Thm.~14]{BLS}。原始族はその部分集合であるから、ここでも2個の加数では十分でない。
本稿の上界と合わせると、
\[
        3\leq h_{\mathrm{asym}}\leq6
\]
を得る。したがって、正確な漸近位数は $3,4,5,6$ のいずれかである。特に、上界6を5へ
下げられるかどうかは未解決である。

\section{OEISの数列}\label{sec:oeis}

本稿では三つのOEIS数列を扱う。完全均衡数は \seqnum{A014486}
\cite{OEIS-A014486} をなす。$\NpD$ の正の元、すなわち原始Dyck数は
\seqnum{A057547} \cite{OEIS-A057547} をなす。本稿で研究する中心的な数列
$a(n)=r(2n)$ は \seqnum{A395858} \cite{OEIS-A395858} である。

\begin{thebibliography}{9}

\bibitem{BHS}
J.~P. Bell, K.~G. Hare, and J.~Shallit,
When is an automatic set an additive basis?,
\emph{Proc. Amer. Math. Soc. Ser. B} \textbf{5} (2018), 50--63.

\bibitem{BLS}
J.~P. Bell, T.~F. Lidbetter, and J.~Shallit,
Additive number theory via approximation by regular languages,
\emph{Internat. J. Found. Comput. Sci.} \textbf{31} (2020), 667--687.

\bibitem{MNRS}
P.~Madhusudan, D.~Nowotka, A.~Rajasekaran, and J.~Shallit,
Lagrange's theorem for binary squares,
in \emph{43rd International Symposium on Mathematical Foundations of
Computer Science}, Leibniz Int. Proc. Inform., Vol.~117, Schloss
Dagstuhl--Leibniz-Zentrum f{\"u}r Informatik, 2018, pp.~18:1--18:14.

\bibitem{Nathanson}
M.~B. Nathanson,
\emph{Additive Number Theory: The Classical Bases},
Grad. Texts in Math., Vol.~164, Springer, 1996.

\bibitem{OEIS-A014486}
OEIS Foundation Inc.,
\emph{The On-Line Encyclopedia of Integer Sequences, A014486},
\url{https://oeis.org/A014486}.

\bibitem{OEIS-A057547}
OEIS Foundation Inc.,
\emph{The On-Line Encyclopedia of Integer Sequences, A057547},
\url{https://oeis.org/A057547}.

\bibitem{OEIS-A395858}
OEIS Foundation Inc.,
\emph{The On-Line Encyclopedia of Integer Sequences, A395858},
\url{https://oeis.org/A395858}.

\bibitem{RSS}
A.~Rajasekaran, J.~Shallit, and T.~Smith,
Sums of palindromes: an approach via automata,
in \emph{35th Symposium on Theoretical Aspects of Computer Science},
Leibniz Int. Proc. Inform., Vol.~96, Schloss Dagstuhl--Leibniz-Zentrum
f{\"u}r Informatik, 2018, pp.~54:1--54:12.

\end{thebibliography}

\begin{flushleft}
Takayuki Kuriyama\\
Growup Learning Academy\\
4-8-3 Shakujii-machi\\
Nerima-ku, Tokyo 177-0041\\
Japan\\
\textit{Email address}: \href{mailto:growup.kuriyama@gmail.com}
{\texttt{growup.kuriyama@gmail.com}}
\end{flushleft}

\end{document}
